% 自旋求和、偏振密度矩阵与外部实光子偏振代数

树级费米子过程反复需要把未分辨外部自旋转成 Dirac 矩阵迹。物理四动量统一取
\begin{equation*}
 p^\mu=(E_{\mathbf p}/c,\mathbf p),
 \qquad p^2=m^2c^2,
\end{equation*}
并使用精细结构常数
\begin{equation*}
 \alpha\equiv\frac{e^2}{4\pi\varepsilon_0\hbar c},
 \qquad e>0,
 \qquad \qe=-e.
\end{equation*}
因此所有自旋子完备关系中的质量项都必须写成动量量纲的 \(mc\)。

正、负能自由旋量的协变完备关系为
\begin{equation}
 \boxed{
 \sum_su_s(p)\bar u_s(p)=\slashed p+mc,
 \qquad
 \sum_sv_s(p)\bar v_s(p)=\slashed p-mc.}
 \label{eq:spin-sum-si-r9}
\end{equation}
结合 Clifford 迹恒等式，矢量流中最常复用的秩二张量直接化为
\begin{equation}
 \boxed{
 L^{\mu\nu}(a,b;m)
 \equiv
 \Tr[(\slashed a+mc)\gamma^\mu(\slashed b+mc)\gamma^\nu]
 =4\!
 \left[a^\mu b^\nu+a^\nu b^\mu
 -(a\!\cdot b-m^2c^2)\eta^{\mu\nu}\right].}
 \label{eq:qed-basic-tensor-si-r9}
\end{equation}
未偏振入射态由自旋密度矩阵 $\rho_s=\mathbb I_2/2$ 定义，末态自旋不分辨则由完备求和实现。将这两个统计操作与\eqnrefs{eq:spin-sum-si-r9} 结合后，矩阵元平方化为相应的 Dirac 迹张量。

若实验准备了确定自旋方向，令 \(\boldsymbol\zeta\) 为粒子静止系中的单位自旋方向，\(|\boldsymbol\zeta|=1\)。把 \((0,\boldsymbol\zeta)\) 与粒子四动量用同一个 Lorentz 推动到实验参考系，得到自旋四矢量
\begin{equation}
 \boxed{
 s^0=\gamma\frac{\mathbf v\!\cdot\boldsymbol\zeta}{c},
 \qquad
 \mathbf s
 =\boldsymbol\zeta
 +\frac{\gamma^2}{\gamma+1}
 \frac{\mathbf v(\mathbf v\!\cdot\boldsymbol\zeta)}{c^2},}
 \label{eq:spin-fourvector-r9}
\end{equation}
满足 \(s^2=-1\) 与 \(s\!\cdot p=0\)。对应的固定自旋密度矩阵为
\begin{equation}
 \boxed{
 u(p,s)\bar u(p,s)
 =\frac12(\slashed p+mc)(1+\gamma^5\slashed s).}
 \label{eq:polarized-spinor-projector-r9}
\end{equation}
对相反的两个自旋方向求和即恢复\eqnrefs{eq:spin-sum-si-r9}。

外部实光子的两个物理自由度已经由\eqnrefs{eq:physical-photon-polarization-projector-dp42} 的横向完备张量定义。若振幅写成 \(\mathcal M=\varepsilon_\mu\mathcal M^\mu\)，并满足 Ward 条件 \(k_\mu\mathcal M^\mu=0\)，则辅助参考矢量相关项全部消失，\eqnrefs{eq:physical-photon-polarization-conserved-current-dp42} 允许在\emph{与守恒振幅完成收缩以后}使用
\begin{equation*}
 \sum_{\lambda=1,2}
 \varepsilon_\lambda^\mu\varepsilon_\lambda^{\nu *}
 \rightsquigarrow-\eta^{\mu\nu}.
\end{equation*}
这个箭头不是完整四维张量恒等式，只表示物理横向偏振求和在守恒矩阵元中的等效替换。

作用量与顶角先确定量子振幅；未分辨的自旋或光子偏振只在振幅构造完成后通过密度矩阵和物理偏振投影进入 $\overline{|\mathcal M|^2}$，再与相空间和入射通量组成实验截面。因而自旋平均、固定自旋和偏振求和属于外态制备与读取规则，不改变相互作用 Hamiltonian。

